\documentclass[12pt]{article}

\usepackage{sbc-template}
\usepackage{graphicx,url}
\usepackage[utf8]{inputenc}
\usepackage[T1]{fontenc}
\usepackage[brazil]{babel}
\usepackage{booktabs}

\sloppy

\title{Provisionamento e Operação de Infraestrutura para Computação\\
Confidencial Aplicada a LLMs: uma Revisão da Literatura}

\author{Raisson Adrian Grangeiro Souto\inst{1}}

\address{Programa de Pós-Graduação em Ciência da Computação (PPGCC)\\
  Universidade Federal de Campina Grande (UFCG) -- Campina Grande, PB -- Brasil
  \email{raisson.souto@copin.ufcg.edu.br}
}

\begin{document}

\maketitle

\begin{abstract}
  % TODO: abstract em inglês (máx. 10 linhas).
\end{abstract}

\begin{resumo}
  % TODO: resumo em português (máx. 10 linhas). Escrever por último.
\end{resumo}

\section{Introdução}\label{sec:intro}
% Contexto: LLMs como serviço multi-tenant expõem pesos, prompts e saídas.
% Problema: TEEs (SGX, TDX, SEV-SNP, GPUs confidenciais) deslocam a
%   complexidade para provisionamento e operação.
% Objetivo e questões de pesquisa (QP1--QP3).
% Contribuições e organização do artigo.

\section{Fundamentação Teórica}\label{sec:fundamentacao}

\subsection{Trusted Execution Environments}
% Enclave (SGX) vs. VM confidencial (TDX, SEV-SNP); unidade de isolamento.

\subsection{Attestation Remota}
% Modelo RATS (Attester, Verifier, Relying Party); liberação de segredos.

\subsection{Orquestração de Cargas Confidenciais}
% Confidential Containers (Kata + Trustee/KBS); ofertas gerenciadas.

\subsection{GPUs Confidenciais e Inferência de LLMs}
% H100/H200 CC mode; extensão da fronteira de confiança.

\section{Metodologia}\label{sec:metodologia}
% Questões de pesquisa; bases (arXiv, ACM DL, IEEE Xplore, USENIX);
% strings de busca; critérios de inclusão/exclusão (2022--2026);
% processo de seleção (n encontrados -> n incluídos); eixos de análise.

\section{Resultados e Discussão}\label{sec:resultados}

\subsection{Eixo 1: Provisionamento de VMs e Containers Confidenciais}
% \cite{misono2024cvms}, \cite{confk8s2025mdpi}, \cite{c8s2026},
% \cite{teecontainers2025trust}

\subsection{Eixo 2: Attestation Remota e Automação de Políticas}
% \cite{narayanan2023evtpm}, \cite{snpguard2024}, \cite{sokattest2025},
% \cite{paju2023sok}, \cite{geppert2022tee}

\subsection{Eixo 3: Computação Confidencial para LLMs e GPUs}
% \cite{chrapek2025confllm}, \cite{zhu2024h100}, \cite{wang2024cpugpu},
% \cite{forough2026agents}, \cite{edgecloud2023ml}

\subsection{Análise Crítica e Lacunas}
% Tabela comparativa dos trabalhos (Tab.~\ref{tab:comparacao}).
% Lacunas: orquestração heterogênea em escala de fleet; CI/CD para builds
% atestáveis; observabilidade sob TEE.

\begin{table}[ht]
\centering
\caption{Comparação dos trabalhos revisados.}\label{tab:comparacao}
\begin{tabular}{@{}llll@{}}
\toprule
Trabalho & Eixo & TEE & Contribuição \\
\midrule
% TODO
\bottomrule
\end{tabular}
\end{table}

\subsection{Ameaças à Validade}
% Viés de seleção, predominância de preprints, recorte temporal.

\section{Conclusão}\label{sec:conclusao}
% Síntese, resposta às QPs, trabalhos futuros (orquestração attestation-aware).

\section*{Uso de Ferramentas de Inteligência Artificial}
% PREENCHER: listar cada ferramenta e para que foi usada neste trabalho.
% Exemplos: busca e triagem de artigos, extração de dados dos estudos,
% revisão de texto, geração de código ou figuras. Remover o que não se aplicar.
Este trabalho utilizou [FERRAMENTA, FORNECEDOR] nas seguintes atividades:
[ATIVIDADES]. Todo o conteúdo gerado com apoio dessas ferramentas foi revisado
e validado pelo autor, e todas as referências foram conferidas nas fontes
originais. O autor assume integral responsabilidade pelo conteúdo, pelas
análises e pelas conclusões.

\bibliographystyle{sbc}
\bibliography{references}

\end{document}
